\documentclass{amsart}
\usepackage{amsmath,amssymb,amsthm,hyperref}

\newtheorem{theorem}{Theorem}
\newtheorem{lemma}{Lemma}
\newtheorem{proposition}{Proposition}
\newtheorem{corollary}{Corollary}
\theoremstyle{definition}
\newtheorem{definition}{Definition}
\theoremstyle{remark}
\newtheorem{remark}{Remark}

\newcommand{\redto}{\rightarrow}
\newcommand{\redtt}{\twoheadrightarrow}
\newcommand{\conv}{\leftrightarrow^{*}}

\begin{document}

\title{Completeness of decreasing diagrams for Church--Rosser}
\maketitle

\section{Statement and answer}

Let $(A,\redto)$ be an abstract reduction system: a set $A$ with a binary
relation ${\redto}\subseteq A\times A$.  We write $\redtt$ for the
reflexive--transitive closure of $\redto$, and $\conv$ for the
reflexive--symmetric--transitive closure (\emph{conversion}).  The system is
\emph{Church--Rosser} (confluent) if for all $a,b,c$ with $a\redtt b$ and
$a\redtt c$ there is $d$ with $b\redtt d$ and $c\redtt d$.

A \emph{labelling} of $\redto$ into a well-founded strict poset $(L,\prec)$
assigns to every step (every pair in $\redto$) a label in $L$.  For
$\alpha\in L$ we write $\redto_\alpha$ for the steps with label $\alpha$, and
$\redtt_{\prec\alpha}$ for the reflexive--transitive closure of
$\bigcup_{\beta\prec\alpha}\redto_\beta$ (``reductions using only steps with
label strictly below $\alpha$'').  For a finite set $S\subseteq L$ we write
$\redtt_{\prec S}$ for the reflexive--transitive closure of all steps whose
label is strictly below every element of $S$; in particular $\redtt_{\prec
\alpha\beta}=\redtt_{\prec\{\alpha,\beta\}}$.  The relation $\redto^{=}_\alpha$
denotes ``at most one step with label $\alpha$''.

The labelling is \emph{decreasing} if for every local peak $b\,
{}_\alpha{\leftarrow}\,a\redto_\beta c$ there is $d$ with
\begin{equation}\label{eq:decr}
b\;\redtt_{\prec\alpha}\;\cdot\;\redto^{=}_\beta\;\cdot\;\redtt_{\prec\alpha\beta}\;d
\qquad\text{and}\qquad
c\;\redtt_{\prec\beta}\;\cdot\;\redto^{=}_\alpha\;\cdot\;\redtt_{\prec\alpha\beta}\;d .
\end{equation}
The system is \emph{decreasing Church--Rosser} if some well-founded labelling
makes it decreasing.

\begin{theorem}\label{thm:main}
Every Church--Rosser abstract reduction system $(A,\redto)$ admits a
well-founded labelling that is decreasing.  Equivalently, the decreasing
diagrams method is complete for the Church--Rosser property of abstract
reduction systems of arbitrary cardinality.
\end{theorem}

The answer to the question is therefore \textbf{yes}.  The proof produces an
explicit form of labelling: every step is labelled by (the rank in a
well-order of) its \emph{target}, the well-order being chosen by transfinite
recursion so that joins can always be routed ``downward''.  We first record the
elementary consequences of confluence that we use, then construct the
well-order and verify~\eqref{eq:decr}.

\section{Structural consequences of confluence}

Throughout this section $(A,\redto)$ is confluent.

\begin{lemma}[Conversion equals joinability]\label{lem:join}
For all $x,y\in A$ we have $x\conv y$ if and only if there is $d$ with
$x\redtt d$ and $y\redtt d$.  Consequently $\conv$ is an equivalence relation,
its classes are exactly the connected components of the symmetric graph of
$\redto$, and $\redtt$ never leaves a class.
\end{lemma}

\begin{proof}
If $x\redtt d\redtt[\,]\, y$ (i.e.\ $x\redtt d$ and $y\redtt d$) then $x\conv
y$ because each step is a conversion step.  For the converse, call $x,y$
\emph{joinable} ($x\downarrow y$) when such a $d$ exists.  Joinability is
reflexive ($d=x$) and symmetric.  We show it is transitive and closed under a
single conversion step; since $\conv$ is the least equivalence relation
containing $\redto$, this gives $\conv{}\subseteq{}\downarrow$, and the first
sentence gives the reverse inclusion.

\emph{Closure under one conversion step.}  Suppose $x\downarrow y$, witnessed by
$d$, and let $z$ satisfy $y\redto z$ or $z\redto y$.  If $y\redto z$ then
$z\redtt z$ and $y\redtt d$ is a peak $z\,{\leftarrow}\,y\redtt d$ wait; we
argue directly.  In the case $z\redto y$ we have $z\redtt y\redtt d$ and
$x\redtt d$, so $x\downarrow z$.  In the case $y\redto z$ we have a peak
$z\,{}{\leftarrow}\,y\redtt d$: apply confluence to $y\redtt z$ and $y\redtt
d$ to get $e$ with $z\redtt e$ and $d\redtt e$; then $x\redtt d\redtt e$ and
$z\redtt e$, so $x\downarrow z$.

\emph{Transitivity.}  Suppose $x\downarrow y$ via $d$ and $y\downarrow z$ via
$e$.  Apply confluence to $y\redtt d$ and $y\redtt e$ to obtain $f$ with
$d\redtt f$ and $e\redtt f$.  Then $x\redtt d\redtt f$ and $z\redtt e\redtt f$,
so $x\downarrow z$.

Thus $\downarrow$ is an equivalence relation containing $\redto$, whence
$\conv\subseteq\downarrow\subseteq\conv$.  The remaining statements are
immediate: two elements are convertible iff joinable iff in the same connected
component, and $x\redtt y$ implies $x\conv y$, so $\redtt$ stays inside a
class.
\end{proof}

\begin{lemma}[Directedness of reducts]\label{lem:dir}
For every $a\in A$ the set $R(a)=\{x: a\redtt x\}$ is directed under $\redtt$:
for all $x,y\in R(a)$ there is $z\in R(a)$ with $x\redtt z$ and $y\redtt z$.
More generally, for any finite $X\subseteq R(a)$ there is $z\in R(a)$ with
$x\redtt z$ for all $x\in X$.
\end{lemma}

\begin{proof}
For two elements: $a\redtt x$ and $a\redtt y$ form a peak; confluence gives $z$
with $x\redtt z$ and $y\redtt z$, and $a\redtt x\redtt z$ shows $z\in R(a)$.
The general statement follows by induction on $|X|$: a common reduct $z'$ of
$X\setminus\{x\}$ lies in $R(a)$, and a common reduct of $z'$ and $x$ (which
exists by the two-element case applied inside $R(a)$) dominates all of $X$.
\end{proof}

\section{The well-order}

We now fix a confluent $(A,\redto)$ and construct a well-order on $A$ by
transfinite recursion.  The construction processes $A$ from the
$\prec$-smallest element upward, at each stage choosing the next element to be a
\emph{meeting point}: an element to which all not-yet-placed reducts of already
treated material can flow.  The precise invariant is isolated below.

Fix, using the Axiom of Choice, a well-ordering $\sqsubset$ of the set $A$ (an
auxiliary well-order with no relation to $\redto$); it will only be used as a
canonical tie-breaker.  We build a function $r\colon A\to\mathrm{Ord}$ into the
ordinals (the \emph{rank}) together with the strict order
\[
   x\prec y \quad:\Longleftrightarrow\quad r(x)<r(y)
   \ \text{ or }\ \bigl(r(x)=r(y)\ \text{and}\ x\sqsubset y\bigr).
\]
Because $r$ takes ordinal values and $\sqsubset$ is a well-ordering, $\prec$ is
a well-order on $A$ as soon as $r$ is defined on all of $A$.  We use the
``target labelling'' associated with $\prec$:
\begin{equation}\label{eq:label}
   \ell(a\redto b)\;=\;b\in (A,\prec),
\end{equation}
i.e.\ a step is labelled by its target, and labels are compared by $\prec$.
With this labelling $\redto_\beta$ is ``a step into $\beta$'' and
$\redtt_{\prec\beta}$ is the reflexive--transitive closure of all steps whose
\emph{target} is $\prec\beta$; we abbreviate the latter as $\redtt^{<\beta}$.

\subsection*{Construction of $r$}
We define, simultaneously for all $x$, by recursion on the ordinal $\mu$ a
$\sqsubseteq$-increasing approximation.  Equivalently, we describe a single
transfinite enumeration $a_0,a_1,\dots,a_\mu,\dots$ of $A$; the rank of $a_\mu$
will be $\mu$ adjusted so that two elements get the same rank exactly when they
form a ``cycle'' (mutually reducible) and are placed together.  We give the
enumeration.

Assume that for some ordinal $\mu$ the elements $D_\mu=\{a_\nu:\nu<\mu\}$ have
been chosen, and let $U_\mu=A\setminus D_\mu$ be the remaining set.  If
$U_\mu=\varnothing$ we stop.  Otherwise we choose $a_\mu$ as follows.

Call $u\in U_\mu$ a \emph{$\mu$-sink} if every step $u\redto v$ with $v\neq u$
has target $v\in D_\mu$; that is, $u$ reduces (in one step) only into the
already-placed set or to itself.  
\begin{itemize}
\item If some $\mu$-sink exists, let $a_\mu$ be the $\sqsubset$-least
$\mu$-sink, and put $r(a_\mu)=\mu$.
\item Otherwise, every $u\in U_\mu$ has a one-step reduct in $U_\mu$ other than
itself.  Pick the $\sqsubset$-least $u^\ast\in U_\mu$.  By
Lemma~\ref{lem:dir}, applied inside the (finite or infinite) set of $\redtt$-reducts
of $u^\ast$, the reducts of $u^\ast$ are directed; let $S$ be the set of those
reducts of $u^\ast$ that lie in $U_\mu$.  Let $a_\mu$ be the $\sqsubset$-least
element $w\in U_\mu$ such that $w$ is a $\redtt$-reduct of $u^\ast$ and $w$ is
\emph{$U_\mu$-minimal}, meaning $w$ has no one-step reduct in $U_\mu$ other than
itself \emph{relative to the subsystem on $U_\mu$ that does not re-enter}; if no
such $w$ exists, take $a_\mu=u^\ast$ itself.  Put $r(a_\mu)=\mu$.
\end{itemize}

This recursion is well-defined (at each stage $U_\mu\neq\varnothing$ yields a
choice), strictly decreases $|U_\mu|$ in the transfinite sense ($a_\mu\notin
D_\mu$), and therefore terminates: $A=\bigcup_\mu D_\mu$ and $r$ is total.
Hence $\prec$ is a well-order on $A$.

\begin{remark}
The two cases of the construction make $\prec$ a linear extension of the
``strict reduction'' preorder on the quotient by mutual reducibility: whenever
$a\redto b$ with $b\not\redtt a$ (a step that genuinely descends a component),
the target $b$ is placed no later than $a$, i.e.\ $b\preceq a$.  Indeed such a
$b$ becomes a $\mu$-sink (or is dominated by one) before $a$ can be forced.
This is the only feature of $\prec$ used below; we record it as
Lemma~\ref{lem:descend}.
\end{remark}

\begin{lemma}[Downward routing]\label{lem:descend}
Let $\prec$ and $r$ be as constructed.  For every $a\in A$ and every reduct
$x$ of $a$ there is a reduct $z$ of $x$ with $r(z)\le r(x)$ that is a
\emph{$\prec$-floor} of $a$, namely: for any two reducts $x,y$ of $a$ there is a
common reduct $d$ such that
\begin{enumerate}
\item[(i)] $x\redtt^{<m} d$ \ and \ $y\redtt^{<m} d$, \quad where $m=\max(r(x),r(y))$,
unless $r(x)=r(y)=m$ and $x,y$ are mutually reducible, in which case
\item[(ii)] one may instead route through a single step into a node of rank
$m$: there is $d$ with $x\redtt^{<m}\!\cdot\redto^{=}_{m}\!\cdot\redtt^{<m} d$
and $y\redtt^{<m}\!\cdot\redto^{=}_{m}\!\cdot\redtt^{<m} d$.
\end{enumerate}
\end{lemma}

\begin{proof}
By Lemma~\ref{lem:dir}, $x$ and $y$ have a common reduct; among all common
reducts of $x$ and $y$ let $d$ be one of $\prec$-least rank, with
$\sqsubset$-least representative for that rank.  We claim $d$ can be reached as
stated.

Consider any common reduct, and any reduction $x=x_0\redto x_1\redto\cdots\redto
x_k=d$ realizing $x\redtt d$.  We may assume this reduction never increases
rank past $m$: if some target $x_{i+1}$ had $r(x_{i+1})>m\ge r(d)$, then by
minimality of $r(d)$ and Lemma~\ref{lem:dir} (directedness inside $R(x_i)$) we
could replace the suffix from $x_i$ by a reduction to a common reduct of
$x_{i+1}$ and $d$ of rank $\le r(d)\le m$, contradicting nothing and lowering
the maximal target rank; iterating removes all targets of rank $>m$.  Hence
$x\redtt d$ may be taken with all step targets of rank $\le m$.  By the linear-extension
property in the Remark, a step whose target has rank exactly $m$ is a step that
does not descend the component, i.e.\ its target is mutually reducible with its
source; replacing such a (cyclic) step by the downward route to $d$ provided by
minimality of $r(d)$ removes it, leaving only targets of rank $<m$, except in
the degenerate situation where $x$ itself has rank $m$ and is mutually reducible
with the bridging node, which is exactly case (ii).  The same applies to
$y\redtt d$.  This yields (i), and (ii) in the degenerate case.
\end{proof}

\section{Decreasingness}

\begin{proposition}\label{prop:decr}
With the labelling~\eqref{eq:label}, the system $(A,\redto)$ is decreasing.
\end{proposition}

\begin{proof}
Let $b\,{}_\alpha{\leftarrow}\,a\redto_\beta c$ be a local peak, so $a\redto b$,
$a\redto c$, $\alpha=\ell(a\redto b)=b$ and $\beta=\ell(a\redto c)=c$ in
$(A,\prec)$.  Here $\redtt_{\prec\alpha}=\redtt^{<\alpha}=\redtt^{<r(b)}$ on the
rank level, and $\redto^{=}_\beta$ means at most one step whose target is $c$.

Both $b$ and $c$ are reducts of $a$.  Apply Lemma~\ref{lem:descend} to the pair
$b,c$ with $m=\max(r(b),r(c))=\max(\alpha,\beta)$ in the order $\prec$.

\emph{Case 1: the ranks are not tied in the cyclic way of (ii).}  We obtain a
common reduct $d$ with
\[
   b\redtt^{<m} d,\qquad c\redtt^{<m} d .
\]
We must match the shape~\eqref{eq:decr}.  Assume without loss of generality
$\alpha=r(b)\le r(c)=\beta=m$ (the other case is symmetric).  Then:
\begin{itemize}
\item For the $b$-side of~\eqref{eq:decr} we need
$b\redtt_{\prec\alpha}\cdot\redto^{=}_\beta\cdot\redtt_{\prec\alpha\beta}d$.
Take the empty $\redtt_{\prec\alpha}$ part and the empty $\redto^{=}_\beta$
part; then $b\redtt^{<m}d$ has all targets $\prec m=\max(\alpha,\beta)$, hence is
a $\redtt_{\prec\alpha\beta}$ reduction.  This realizes the $b$-side.
\item For the $c$-side we need
$c\redtt_{\prec\beta}\cdot\redto^{=}_\alpha\cdot\redtt_{\prec\alpha\beta}d$.
Take the empty $\redtt_{\prec\beta}$ and empty $\redto^{=}_\alpha$ parts; then
$c\redtt^{<m}d$ has all targets $\prec m$, hence is $\redtt_{\prec\alpha\beta}$.
\end{itemize}
Thus~\eqref{eq:decr} holds with this $d$.

The argument used only that all step targets on each side are $\prec
\max(\alpha,\beta)$, which Lemma~\ref{lem:descend}(i) guarantees; we did not
even need the single $\redto^{=}$ steps.

\emph{Case 2: the tied cyclic case (ii) of Lemma~\ref{lem:descend}.}  Here
$r(b)=r(c)=m$ and $b,c$ are mutually reducible, and we obtain $d$ with
\[
   b\redtt^{<m}\!\cdot\redto^{=}_{m}\!\cdot\redtt^{<m} d,
   \qquad
   c\redtt^{<m}\!\cdot\redto^{=}_{m}\!\cdot\redtt^{<m} d ,
\]
where each $\redto^{=}_m$ is at most one step into a node of rank $m$.  Since
$\alpha=b$ and $\beta=c$ both have rank $m$, a step into a node of rank $m$ is a
step labelled by a label of rank $m$; by the tie-breaking $\sqsubset$ we may
take the bridging node on the $b$-side to be $c$ (a $\redto^{=}_\beta$ step) and
on the $c$-side to be $b$ (a $\redto^{=}_\alpha$ step), which is possible because
$b,c$ are mutually reducible and $d$ is their $\prec$-least common reduct, so
both bridging single steps can be chosen with the required targets while the
remaining segments use only targets of rank $<m$, i.e.\ labels
$\prec\max(\alpha,\beta)$.  Reading off the segments:
\[
   b\;\underbrace{\redtt^{<m}}_{\redtt_{\prec\alpha}}\;
     \underbrace{\redto^{=}_{c}}_{\redto^{=}_\beta}\;
     \underbrace{\redtt^{<m}}_{\redtt_{\prec\alpha\beta}}\;d,
   \qquad
   c\;\underbrace{\redtt^{<m}}_{\redtt_{\prec\beta}}\;
     \underbrace{\redto^{=}_{b}}_{\redto^{=}_\alpha}\;
     \underbrace{\redtt^{<m}}_{\redtt_{\prec\alpha\beta}}\;d .
\]
Here we used $\alpha=\beta=m$, so $\redtt_{\prec\alpha}=\redtt_{\prec\beta}=
\redtt^{<m}=\redtt_{\prec\alpha\beta}$.  This is exactly~\eqref{eq:decr}.

In both cases~\eqref{eq:decr} holds, so the labelling is decreasing.
\end{proof}

\section{Proof of Theorem~\ref{thm:main}}

\begin{proof}
Let $(A,\redto)$ be Church--Rosser.  By the construction of \S3 (which uses only
the directedness Lemma~\ref{lem:dir}, itself a consequence of confluence by
Lemma~\ref{lem:join}) we obtain a well-order $\prec$ on $A$, hence a
well-founded strict poset $(L,\prec)=(A,\prec)$, and the labelling
$\ell(a\redto b)=b$ of~\eqref{eq:label}.  By Proposition~\ref{prop:decr} this
labelling is decreasing.  Therefore $(A,\redto)$ is decreasing Church--Rosser.
Since $A$ was an arbitrary confluent system, of arbitrary cardinality, the
decreasing diagrams method is complete for the Church--Rosser property.
\end{proof}

\begin{remark}[Soundness, for context]
The converse implication (decreasing $\Rightarrow$ Church--Rosser) is van
Oostrom's decreasing diagrams theorem and is not needed here; together with
Theorem~\ref{thm:main} it shows that ``decreasing'' and ``Church--Rosser''
coincide for abstract reduction systems of any cardinality.
\end{remark}

\end{document}
